\chapter{Theoretical Background}
\label{ch:theory}

This chapter derives, in full, the theory whose motivation and literature
context were reviewed qualitatively in \cref{ch:litreview}. It proceeds in
two parts. \Cref{sec:th-migration} derives the three migration algorithms
used to focus the raw GPR data
(\cref{ch:resolution,ch:tl-horizontal,ch:tl-vertical,ch:fluidflow}).
\Cref{sec:th-fourier-shift}--\ref{sec:th-local} then derive the 2D phase-plane
shift-estimation method that is validated in \cref{ch:phase-plane} and applied
to a realistic scenario in \cref{ch:fluidflow}: starting from the Fourier shift
theorem, building the weighted least-squares (WLS) plane fit, showing exactly
why it decouples geometric movement from material change, and finally
connecting the global (wavenumber-domain) and local (trace-by-trace,
time-frequency) views of the same physics.

\section{GPR Migration Fundamentals}
\label{sec:th-migration}

All data in this thesis are zero-offset (collocated transmitter/receiver)
B-scans, which makes the \emph{exploding-reflector model} applicable
\citep{claerbout1985}: every reflector in the subsurface is treated as if it
were an active source that radiates a pulse upward at $t=0$, recorded by
receivers at the surface. Because the real two-way travel time corresponds to
a wave travelling down and back up once, this fictitious one-way exploding
source must propagate at half the true medium velocity,
\begin{equation}
  \vmig = \frac{v}{2},
  \label{eq:vmig}
\end{equation}
so that the one-way travel time in the exploding-reflector model exactly
matches the two-way travel time of the real survey. Every migration algorithm
used in this thesis operates under this convention.

\subsection{Kirchhoff (Delay-and-Sum) Migration}

Kirchhoff migration is formulated as a linear forward operator $K$ mapping a
reflectivity model $m(z,x)$ to recorded data $d(t,x)$ by summing the
reflectivity along the travel-time hyperbola of every trace,
\begin{equation}
  t(x, x_s, z) = \frac{r}{\vmig} + t_0 ,\qquad r = \sqrt{z^2 + (x-x_s)^2},
  \label{eq:kirchhoff-traveltime}
\end{equation}
where $x_s$ is the source/receiver position and $t_0$ a static time-zero
correction. Migration is the adjoint operation: every recorded sample is
smeared back (\emph{delay-and-sum}) along its hyperbola of possible
origins, and the image is the sum over all traces,
\begin{equation}
  m_{\mathrm{mig}} = K^{\mathsf{T}} d .
  \label{eq:kirchhoff-adjoint}
\end{equation}
This thesis uses the PyLops zero-offset Kirchhoff operator with an
analytic-signal Ricker wavelet matched to the source, evaluated with a finite
migration aperture \citep{schneider1978}.

\subsection{Gazdag Phase-Shift Migration}

Gazdag migration works entirely in the frequency--wavenumber ($f$-$k_x$)
domain \citep{gazdag1978}. The recorded wavefield is downward-continued one
depth step $\delta z$ at a time by multiplying its 2D temporal-frequency /
horizontal-wavenumber spectrum by a phase-shift operator,
\begin{equation}
  U(z+\delta z, k_x, \omega) = U(z, k_x,\omega)\, e^{j k_z \delta z},
  \qquad
  k_z = \sqrt{\left(\frac{\omega}{\vmig}\right)^2 - k_x^2},
  \label{eq:gazdag}
\end{equation}
and the image is built by applying the imaging condition --- extracting the
$t=0$ component of the continued field --- at every depth step. Bins for
which the argument of the square root in \cref{eq:gazdag} is negative
correspond to evanescent energy and are set to zero before continuation;
otherwise the unstable exponential growth of an imaginary $k_z$ produces
migration "smile" artefacts. This thesis pads each B-scan with a $5\,\%$
cosine taper and $100\,\%$ zero-padding in $x$ before transforming, to
suppress wrap-around.

\subsection{Back-Propagation (Time-Reversal) Migration}

As an independent, purely numerical cross-check of the two analytic methods
above, every B-scan is also migrated by literal time-reversal: each trace is
reversed in time, normalised, and re-injected as a source at its original
receiver position into a finite-difference time-domain (gprMax) model of a
homogeneous medium at $\vmig$. By the time-reversal symmetry of the wave
equation, the back-propagated field refocuses at the true scatterer location
at the focusing time
\begin{equation}
  t_{\mathrm{focus}} = T - t_0 ,
  \label{eq:backprop-focus}
\end{equation}
where $T$ is the trace length. The migrated image is read off as the field
snapshot at $t_{\mathrm{focus}}$, either as the full electric-field magnitude
$\lVert\mathbf{E}\rVert$ or as the single polarised component $E_z$. Unlike
Kirchhoff and Gazdag migration, this method makes no high-frequency or
zero-offset approximation beyond the exploding-reflector velocity halving
itself, which makes it a useful independent check on the other two.

\section{The 2D Fourier Shift Theorem}
\label{sec:th-fourier-shift}

Let the baseline migrated image be $b(z,x)$. If a point scatterer
translates by a vertical distance $\Dz$ and a lateral distance $\Dx$ between
the baseline and monitor survey, the monitor image is, to the extent that
migration is linear and the two surveys are migrated with the same velocity
model, a perfectly translated copy of the baseline image,
\begin{equation}
  m(z,x) = b(z - \Dz,\, x - \Dx).
  \label{eq:spatial-shift}
\end{equation}
Taking the 2D continuous Fourier transform of the baseline image,
\begin{equation}
  B(\kz,\kx) = \int_{-\infty}^{\infty}\!\int_{-\infty}^{\infty}
    b(z,x)\, e^{-j(\kz z + \kx x)}\, dz\, dx,
  \label{eq:2d-fourier}
\end{equation}
the Fourier shift theorem states that the spatial translation in
\cref{eq:spatial-shift} becomes a linear phase rotation of the spectrum,
\begin{equation}
  M(\kz,\kx) = B(\kz,\kx)\cdot e^{-j(\kz \Dz + \kx \Dx)}.
  \label{eq:fourier-shift-theorem}
\end{equation}
The amplitude spectrum is unchanged, $|M|=|B|$, but every frequency
coordinate $(\kz,\kx)$ acquires a phase shift proportional to that
coordinate. This is the central fact the rest of the chapter exploits: a
displacement that is invisible in the migrated amplitude image
(\cref{ch:resolution}) is, in principle, exactly recoverable from the phase of
the spectrum.

\section{Isolating the Phase Plane: The Cross-Spectrum}
\label{sec:th-cross-spectrum}

The absolute phase of a single migrated image is not directly useful: it
depends on the shape of the source wavelet and on the (arbitrary) position of
the target within the image, and is, in general, a chaotic, wrapped function
of $(\kz,\kx)$. The displacement information in \cref{eq:fourier-shift-theorem}
is isolated by forming the complex \emph{cross-spectrum} between the baseline
and monitor spectra,
\begin{equation}
  \XS(\kz,\kx) = B(\kz,\kx)\cdot M^{*}(\kz,\kx).
  \label{eq:cross-spectrum-def}
\end{equation}
Substituting \cref{eq:fourier-shift-theorem},
\begin{equation}
  \XS(\kz,\kx) = B(\kz,\kx)\cdot\bigl[B(\kz,\kx)\, e^{-j(\kz \Dz + \kx \Dx)}\bigr]^{*}
  = |B(\kz,\kx)|^{2}\, e^{j(\kz \Dz + \kx \Dx)} .
  \label{eq:cross-spectrum-expand}
\end{equation}
Because $|B|^2$ is real and positive, extracting the phase of $\XS$ via
$\Phi = \angle \XS$ makes it vanish completely, leaving
\begin{equation}
  \boxed{\Phi(\kz,\kx) = \kz \Dz + \kx \Dx} ,
  \label{eq:phase-plane}
\end{equation}
a perfectly flat plane through the origin: its slope along $\kz$ is exactly
the vertical displacement $\Dz$, and its slope along $\kx$ is exactly the
lateral displacement $\Dx$.

\paragraph{Units of the cross-phase spectrum.} The value of $\Phi$ is an
angle (rad or deg); its domain, the wavenumbers $(\kz,\kx)$, has units of
$\mathrm{rad\,m^{-1}}$. The slope of the plane in \cref{eq:phase-plane},
$\partial\Phi/\partial\kz$, therefore has units of
$\mathrm{rad} / (\mathrm{rad\,m^{-1}}) = \mathrm{m}$: the radians cancel, and
what remains is a physical distance. This is the precise sense in which an
angular rotation in the frequency domain converts directly into a
millimetre-scale displacement.

\section{Weighted Least-Squares Plane Fitting}
\label{sec:th-wls}

A discretised GPR image yields not one but $N$ frequency-bin observations of
\cref{eq:phase-plane} inside a chosen passband (\cref{sec:th-mask-weight}),
each of the form
\begin{equation}
  \Phi_i = \kz_{,i}\,\Dz + \kx_{,i}\,\Dx + c, \qquad i=1,\dots,N,
  \label{eq:phase-plane-discrete}
\end{equation}
where $c$ is a small constant absorbing calibration bias (and, as shown in
\cref{sec:th-material-change}, the signature of material change). Collecting
all $N$ equations into the overdetermined linear system $A\mathbf{u}=\bm{\Phi}$,
\begin{equation}
  \underbrace{\begin{bmatrix}
    \kz_{,1} & \kx_{,1} & 1 \\
    \kz_{,2} & \kx_{,2} & 1 \\
    \vdots   & \vdots   & \vdots \\
    \kz_{,N} & \kx_{,N} & 1
  \end{bmatrix}}_{A}
  \underbrace{\begin{bmatrix} \Dz \\ \Dx \\ c \end{bmatrix}}_{\mathbf{u}}
  =
  \underbrace{\begin{bmatrix} \Phi_1 \\ \Phi_2 \\ \vdots \\ \Phi_N \end{bmatrix}}_{\bm{\Phi}}.
  \label{eq:design-matrix}
\end{equation}

\subsection{Why ordinary least squares fails}

Ordinary least squares (OLS) minimises $S(\mathbf{u}) = (\bm\Phi - A\mathbf{u})^{\mathsf T}(\bm\Phi - A\mathbf{u})$,
giving the closed-form solution $\mathbf{u} = (A^{\mathsf T}A)^{-1}A^{\mathsf T}\bm\Phi$. This
treats every frequency bin as equally reliable, but in a GPR spectrum a bin
at the antenna's peak power is far more reliable than a bin at the edge of
the band, which is dominated by background noise. OLS gives both bins equal
weight, letting noisy bins corrupt the plane fit.

\subsection{The weighted solution}
\label{sec:th-mask-weight}

Two safeguards make the fit robust enough for sub-wavelength accuracy.

\paragraph{A. The band-pass mask.} The system in \cref{eq:design-matrix} is
restricted to bins inside the coherent envelope of the source wavelet,
$|\kz|,|\kx| < 1.4\,k_{zc}$, where $k_{zc}$ is the dominant vertical
wavenumber of the pulse. This keeps the total phase rotation $\Phi_i$ inside
$\pm\pi$, preventing the fit from wrapping.

\paragraph{B. Amplitude weighting.} A diagonal weight matrix $W$ is built from
the cross-spectrum magnitude, $W_{ii} = |\XS_i|$, so that high-energy bins
dominate the fit and noise-floor bins are suppressed. The cost function
becomes the energy-weighted residual,
\begin{equation}
  S_W(\mathbf{u}) = \sum_{i=1}^{N} |\XS_i|^{2}
    \Bigl(\Phi_i - (\kz_{,i}\Dz + \kx_{,i}\Dx + c)\Bigr)^{2},
  \label{eq:wls-cost}
\end{equation}
which is minimised by
\begin{equation}
  \mathbf{u} = \bigl(A^{\mathsf T}W^{2}A\bigr)^{-1}A^{\mathsf T}W^{2}\bm\Phi .
  \label{eq:wls-solution}
\end{equation}
In practice, computing $(A^{\mathsf T}W^2 A)^{-1}$ directly squares the
condition number of the system; the implementation instead pre-multiplies
both sides by $W$,
\begin{equation}
  (WA)\,\mathbf{u} = (W\bm\Phi),
  \label{eq:wls-preweighted}
\end{equation}
and solves the pre-weighted system with \texttt{numpy.linalg.lstsq} via
singular value decomposition, which is equivalent to \cref{eq:wls-solution}
but numerically far more stable.

\section{Decoupling Geometric Movement from Material Change}
\label{sec:th-material-change}

The third column of the design matrix in \cref{eq:design-matrix} is the
constant vector $\mathbf{1}$, deliberately included alongside the two
wavenumber columns. This section shows why that single design choice lets the
same fit simultaneously measure displacement \emph{and} detect a change in
the dielectric material at the target, and why the two never contaminate
each other.

\subsection{Why a sub-wavelength fracture produces a frequency-independent phase shift}

Consider a fracture of thickness $d$ much smaller than the wavelength,
$d\ll\lambda$, embedded in a host matrix of permittivity $\eps_m$ and filled
with a material of permittivity $\eps_f$. The wave reflects once at the top
and once at the bottom of the fracture; for $d \ll \lambda$ these two
reflections overlap into a single composite wavelet, and the total
reflection coefficient can be approximated as
\begin{equation}
  R_{\mathrm{total}} \approx R_{\mathrm{top}} + R_{\mathrm{bottom}}\, e^{-j2kd}.
  \label{eq:thinlayer-1}
\end{equation}
A first-order expansion of the exponential for $kd \ll 1$ gives
\begin{equation}
  R_{\mathrm{total}} \approx -j\omega\, \frac{d\sqrt{\eps_m}}{2c_0}
    \left(\frac{\eps_m - \eps_f}{\eps_f}\right),
  \label{eq:thinlayer-2}
\end{equation}
where $c_0$ is the speed of light in vacuum. The leading $-j\omega$ factor
means a thin layer reflects the time-derivative of the incident pulse.

When the fracture fluid changes between the baseline and monitor survey
(e.g.\ air, $\eps_{\mathrm{air}}=1$, displaced by water, $\eps_{\mathrm{water}}=81$)
while the thickness $d$ stays fixed, the ratio of the two reflection
coefficients is
\begin{equation}
  \Delta R = \frac{R_{\mathrm{mon}}}{R_{\mathrm{base}}}
  = \frac{\eps_{\mathrm{air}}\,(\eps_m - \eps_{\mathrm{water}})}
         {\eps_{\mathrm{water}}\,(\eps_m - \eps_{\mathrm{air}})} .
  \label{eq:thinlayer-ratio}
\end{equation}
Both the geometric factor $d$ and the derivative factor $-j\omega$ cancel
exactly in this ratio, leaving a complex constant, independent of $\omega$
and therefore independent of $\kz$ and $\kx$. In the phase domain this is a
uniform rotation
\begin{equation}
  \Delta R \approx \alpha\, e^{-j\Dtheta},
  \label{eq:thinlayer-phase}
\end{equation}
applied equally to every frequency in the pulse bandwidth: a flat,
frequency-independent phase offset $\Dtheta$, with an amplitude attenuation
$\alpha$.

\subsection{Why the WLS fit puts material change exactly into the intercept}

For a fluid-substitution event with no mechanical movement
($\Dz=\Dx=0$), \cref{eq:thinlayer-phase} means every observation in
\cref{eq:phase-plane-discrete} is the same constant, $\Phi_i = \Dtheta$ for
all $i$. Looking at the three columns of $A$ in \cref{eq:design-matrix}:
\begin{itemize}
  \item the $\kz$ column spans negative to positive wavenumbers --- any
    non-zero $\Dz$ tilts the predicted plane, which can only \emph{increase}
    the residual against a perfectly flat target, so the optimum is $\Dz=0$;
  \item the $\kx$ column behaves identically, forcing $\Dx=0$;
  \item the constant column is exactly $\mathbf 1$, so setting $c=\Dtheta$
    matches the flat target with zero residual.
\end{itemize}
This is possible because the wavenumber columns and the constant column are
linearly independent (orthogonal in the sense relevant to least squares): the
matrix inversion in \cref{eq:wls-solution} decouples them exactly, with no
approximation. If a fracture both moves \emph{and} fills with fluid at once,
the fitted plane both tilts (giving $\Dz,\Dx$) and shifts vertically (giving
$c=\Dtheta$), and the two effects remain perfectly separable in the same
single fit, with the intercept tracking the dielectric contrast,
\begin{equation}
  c = \Dtheta \approx \frac{2\pi d}{\lambda}
    \left(\sqrt{\eps_{\mathrm{fluid}}} - \sqrt{\eps_{\mathrm{baseline}}}\right).
  \label{eq:intercept-material}
\end{equation}
This orthogonality is exploited directly in \cref{ch:fluidflow}, where the
target is not a discrete scatterer but a spatially extended fluid front.

\section{Space-Wavenumber Duality: Lateral versus Vertical Asymmetry}
\label{sec:th-duality}

The phase-plane fit of \cref{sec:th-wls} is a \emph{global} operation: it
acts on the whole 2D spectrum at once. \Cref{ch:phase-plane} also examines a
\emph{local} alternative --- the spatial instantaneous phase, extracted
trace-by-trace with a Riesz/Hilbert transform. This section shows that the
two views are consistent, and that they reveal an important physical
asymmetry between the lateral and vertical directions of a migrated image.

\subsection{Lateral direction: position as a proxy for wavenumber}

Near the apex of a migrated point scatterer, the lateral instantaneous phase
$\phi(x)$ is well approximated by a parabola, $\phi(x)\approx\tfrac12 Cx^2+\phi_0$,
where $C$ is the spatial curvature of the focused pulse. Its spatial
derivative, the instantaneous lateral wavenumber, is then
\begin{equation}
  k_{x,\mathrm{inst}}(x) = \frac{\partial\phi}{\partial x} \approx C\,x,
  \label{eq:kx-inst}
\end{equation}
i.e.\ \emph{linearly proportional to lateral position}: the migrated image
behaves like a prism, naturally separating horizontal spatial frequencies
across its width. If the scatterer shifts laterally by $\Dx$, a first-order
Taylor expansion of $\phi_2(x)=\phi_1(x-\Dx)$ gives a phase-difference image
\begin{equation}
  \Delta\phi(x) = \phi_2(x)-\phi_1(x) \approx -k_{x,\mathrm{inst}}(x)\,\Dx
    = -C\,x\,\Dx ,
  \label{eq:dphi-lateral}
\end{equation}
i.e.\ a straight \emph{slope} in $x$, whose steepness is proportional to
$\Dx$ — consistent with \cref{eq:phase-plane} evaluated along the $\kx$ axis.

\subsection{Vertical direction: a constant plateau, not a slope}

The vertical trace through a migrated scatterer is, by contrast, a
compressed source pulse oscillating at the dominant vertical carrier
wavenumber $k_{zc}=4\pi f_c/v$. Near the peak, its phase is linear in depth,
$\phi(z)\approx k_{zc}(z-z_0)+\phi_0$, so the instantaneous vertical
wavenumber is
\begin{equation}
  k_{z,\mathrm{inst}}(z) = \frac{\partial\phi}{\partial z} \approx k_{zc}
  = \text{constant} ,
  \label{eq:kz-inst}
\end{equation}
independent of $z$: unlike the lateral case, depth position is \emph{not} a
proxy for vertical wavenumber, because the vertical trace does not spatially
separate its frequency content the way the migrated lateral profile does.
A vertical displacement $\Dz$ then produces
\begin{equation}
  \Delta\phi(z) \approx -k_{zc}\,\Dz ,
  \label{eq:dphi-vertical}
\end{equation}
which contains no remaining $z$ dependence: $\partial[\Delta\phi(z)]/\partial z\approx0$.
A vertical shift therefore appears in the spatial domain as a flat
\emph{plateau} whose constant level is proportional to $\Dz$, not as a slope.

\subsection{Reconciling the two pictures}

Combining \cref{eq:dphi-lateral,eq:dphi-vertical}, the full 2D spatial phase
difference near the apex $(z_0,0)$ of a scatterer displaced by
$(\Dz,\Dx)$ is
\begin{equation}
  \Delta\phi(z,x) \approx -\bigl(k_{zc}\,\Dz + (C\,x)\,\Dx\bigr),
  \label{eq:dphi-2d}
\end{equation}
which has exactly the same structure as the global phase plane in
\cref{eq:phase-plane}, with the local approximations $k_{z,\mathrm{inst}}\to\kz$
and $C\,x \to \kx$ replacing the true global wavenumber coordinates. The
wavenumber-domain fit of \cref{sec:th-wls} is preferred over this spatial
picture for three reasons: (i) the parabolic approximation in
\cref{eq:kx-inst} only holds near the apex, whereas the Fourier plane in
\cref{eq:phase-plane} is exactly flat everywhere inside the passband; (ii)
amplitude weighting by $|\XS|$ has no clean spatial-domain analogue, since the
spatial phase-difference image does not distinguish signal from background
noise; and (iii) in the spatial picture a material-change intercept
$\Dtheta$ mixes irrecoverably with the vertical term $-k_{zc}\Dz$ in
\cref{eq:dphi-2d}, whereas the wavenumber-domain fit keeps them exactly
orthogonal (\cref{sec:th-material-change}).

\section{Time-Frequency Perspective: Local Phase and Spectral-Line Analysis}
\label{sec:th-local}

The phase-plane fit treats the migrated image as a whole. \Cref{ch:phase-plane}
(\cref{sec:tlp-spectral-line,sec:tlp-spectrogram,sec:tlp-stft}) instead
analyses individual unmigrated or migrated \emph{traces} with a localised
time-frequency transform, which gives access to \emph{when} (in two-way
time) a phase change occurs, complementing the spatial picture above.

\subsection{The localised Fourier shift theorem}

Let $s_1(t)$ be a baseline trace. A sliding-window transform (Gaussian or
Hanning window $w$) gives a localised spectrum
\begin{equation}
  S_1(\tau,\omega) = \int_{-\infty}^{\infty} s_1(t)\,w(t-\tau)\,e^{-j\omega t}\,dt ,
  \label{eq:stft-def}
\end{equation}
at window centre $\tau$ and angular frequency $\omega=2\pi f$. If the monitor
trace is a delayed, phase-rotated copy, $s_2(t) = s_1(t-\Dt)\,e^{j\Dtheta}$,
and the delay $\Dt$ is small compared with the window width, the window
itself barely shifts and the \emph{localised Fourier shift theorem} applies,
\begin{equation}
  S_2(\tau,\omega) \approx S_1(\tau,\omega)\, e^{-j\omega\Dt}\, e^{j\Dtheta} .
  \label{eq:local-shift}
\end{equation}
Forming the local cross-spectrum exactly as in \cref{eq:cross-spectrum-def}
and taking its angle, the baseline phase and amplitude cancel, leaving the
local analogue of \cref{eq:phase-plane},
\begin{equation}
  \boxed{\Delta\Phi(\tau,f) \approx -2\pi f\,\Dt + \Dtheta} .
  \label{eq:local-phase-line}
\end{equation}
At the two-way time $\tau_0$ of the target reflection, a straight-line fit of
$\Delta\Phi$ against $f$ has slope $-2\pi\Dt$ (the mechanical shift) and
intercept $\Dtheta$ (the material change) — the exact time-domain counterpart
of the slope/intercept decomposition in \cref{sec:th-wls}. Converting between
the two-way-time and depth pictures uses the standard relation
\begin{equation}
  \Dt = \frac{2\Dz}{v} ,
  \label{eq:dt-dz}
\end{equation}
so that $\omega\Dt \equiv \kz\Dz$ with $\kz = 2\omega/v$: the temporal-frequency
slope and the vertical-wavenumber slope are the same physical quantity, viewed
in two different but exactly equivalent coordinate systems.

\subsection{Three diagnostic views used in Chapter~\ref{ch:phase-plane}}

Three derived plots make \cref{eq:local-phase-line} directly visible in the
data, and are used repeatedly in the figures of \cref{ch:phase-plane}:

\begin{description}
  \item[Spectral line, $\Delta\Phi(f)$ at fixed $\tau_0$] --- a straight line
    through the origin for pure mechanical movement; a flat line offset from
    zero for pure material change; a sloped line with non-zero intercept for
    a combination of the two.
  \item[Cross-phase spectrogram, $\Delta\Phi(\tau,f)$] --- outside the target
    reflection this is incoherent, salt-and-pepper phase noise; at the
    target's two-way time a coherent window appears, showing a vertical
    fringe pattern for movement or a uniform colour block for a pure material
    change.
  \item[Polar vector rotation] --- the complex coefficient at the dominant
    frequency and peak two-way time, plotted as a vector in the complex
    plane for baseline and monitor; a material change rotates this vector by
    exactly $\Dtheta$, with negligible length change for a purely geometric
    shift.
\end{description}

These local, trace-based views and the global 2D wavenumber fit of
\cref{sec:th-wls} are not competing techniques: they are Fourier duals of the
same underlying physics, related by a spatial Fourier transform of the
time-frequency decomposition with respect to the lateral coordinate $x$,
which reduces (under a smooth-window approximation) directly to the global
spectrum $U(\omega,\kx)$ used throughout this chapter. The wavenumber-domain
fit remains the primary quantitative tool because it linearises the spatial
curvature of \cref{eq:kx-inst} exactly and admits amplitude weighting
natively (\cref{sec:th-mask-weight}); the local, time-frequency view is used
in \cref{ch:phase-plane} to localise \emph{where} (in $\tau$) a phase
anomaly originates, which the global fit alone cannot show.
